% from ega4-9-fr.tex
\oldpage[IV-3]{55}
Soit alors $\Omega$ une extension algébriquement close de $k$ (dans l'univers
considéré). Les conditions suivantes sont équivalentes~:
\begin{enumerate}[label=\emph{\alph*)}]
  \item $\mathcal C$ est non vide.
  \item Il existe $K'\in\mathcal C$ qui est une extension finie de $k$.
  \item On a $\Omega\in\mathcal C$.
\end{enumerate}
\end{proposition}

La condition (i) implique évidemment que \emph{b)} entraîne \emph{c)}, et
\emph{c)} entraîne trivialement \emph{a)}~; prouvons que \emph{a)} entraîne
\emph{b)}. En vertu de (ii) et (iii) il existe une extension
$K'\in\mathcal C$ qui est de la forme $A/\mathfrak m$, où $A$ est une
$k$-algèbre de type fini sur $k$ et $\mathfrak m$ un idéal maximal de $A$.
On sait, par le théorème des zéros de Hilbert (Bourbaki, \emph{Alg. comm.},
chap.~V, \S~3, n\textsuperscript{o}~1, cor.~2 du th.~1) que $K'$ est une
extension finie de $k$.

\begin{corollary}[9.1.2]
\label{IV.9.1.2-fr}
Sous les hypothèses (i), (ii), (iii) de (9.1.1), si $k$ est algébriquement
clos et si $\mathcal C$ n'est pas vide, $\mathcal C$ contient toutes les
extensions de $k$ dans l'univers envisagé.
\end{corollary}

En effet, comme \emph{a)} entraîne \emph{c)}, on a $k\in\mathcal C$, et la
conclusion résulte de l'hypothèse (i).

\begin{remark}[9.1.3]
\label{IV.9.1.3-fr}
En pratique, on vérifiera la condition (ii) de (9.1.1) en notant que $K$ est
limite inductive de ses sous-extensions de type fini $K_\alpha$, et en
appliquant les résultats du \S~8, compte tenu éventuellement du fait que pour
$K_\alpha\subset K_\beta$, $K_\beta$ est fidèlement plat sur $K_\alpha$.
Fréquemment l'ensemble $\mathcal C$ est formé des \emph{corps} appartenant à
un ensemble $\mathcal C'$ de $k$-algèbres qui vérifie la condition suivante~:
\begin{enumerate}[label=(i bis)]
  \item Si $A\in\mathcal C'$ et s'il existe un homomorphisme de $k$-algèbres
  $A\to A'$ (où $A'$ appartient à l'univers où l'on se place), alors
  $A'\in\mathcal C'$.
\end{enumerate}
Lorsqu'il en est ainsi, la condition (i) est trivialement satisfaite, et on
vérifiera généralement la condition (iii) de (9.1.1) en notant que le corps
des fractions $K$ de $A$ est la limite inductive des algèbres $A_f$ (pour la
relation $D(f)\supset D(g)$), et en appliquant les résultats du \S~8, compte
tenu éventuellement du fait que les morphismes $D(g)\to D(f)$ sont des
immersions ouvertes.

Par contre, lorsque (i bis) n'est pas vérifiée, la démonstration de (iii) est
souvent plus délicate, et est liée à des critères de constructibilité qui
seront développés plus loin.
\end{remark}

Voici des exemples typiques d'application du principe de l'extension finie~:

\begin{proposition}[9.1.4]
\label{IV.9.1.4-fr}
Soient $k$ un corps, $\Omega$ une extension algébriquement close de $k$,
$X$ et $Y$ deux préschémas de type fini sur $k$. Les conditions suivantes
sont équivalentes~:
\begin{enumerate}[label=\emph{\alph*)}]
  \item Il existe un $\Omega$-morphisme $X_{(\Omega)}\to Y_{(\Omega)}$
  (resp. un $\Omega$-morphisme possédant l'une (déterminée) des propriétés
  (i) à (xiv) de (8.10.5)).
  \item Il existe une extension finie $k'$ de $k$ et un $k'$-morphisme
  $X_{(k')}\to Y_{(k')}$ (resp. un $k'$-morphisme possédant la propriété
  considérée).
  \item Il existe une extension $K$ de $k$ et un $K$-morphisme
  $X_{(K)}\to Y_{(K)}$ (resp. un $K$-morphisme possédant la propriété
  considérée).
\end{enumerate}
\end{proposition}

On applique la remarque (9.1.3) en prenant pour $\mathcal C'$ l'ensemble de
toutes les $k$-algèbres $A$ (de l'univers où on est placé) telles qu'il existe
un $A$-morphisme $X\otimes_k A\to Y\otimes_k A$ (resp. un morphisme ayant
celle des propriétés de (8.10.5) que l'on considère). La condition (i bis) de
(9.1.3) est alors vérifiée grâce au fait que les

